\begin{name}
	{\tenchude}
	{TOÁN 10}
	{LỚP TOÁN THẦY PHÁT}
	{{Ai cũng sẽ phải trưởng thành dù muốn hay không.\\ Vậy hãy chọn cách mạnh mẽ mà trưởng thành.}}
\end{name}
\Opensolutionfile{ansbook}[ans/ansbookDe6]
\TN
\Opensolutionfile{ans}[ans/ansDe1-TN6]
\begin{ex}%[10-11-EX-GHK2-24-25]%[Dự Đỗ]%[0D3H1-5]
Cho hàm số $y=f(x)$ có đồ thị như hình vẽ.
\begin{center}
\begin{tikzpicture}[line join=round, line cap=round,>=stealth]
\tikzset{every node/.style={scale=0.9}}
\draw[->] (-2.1,0)--(4,0) node[below left] {$x$};
\draw[->] (0,-4.1)--(0,2.1) node[below left] {$y$};
\draw (0,0) node [below left] {$O$};
\foreach \x/\nx in {2/2}
\draw[] (\x,1pt)--(\x,-1pt) node [above] {$\nx$};
\foreach \y/\ny in {-3/-3,1/1}
\draw[] (1pt,\y)--(-1pt,\y) node [above left] {$\ny$};
\draw[dashed](2,0)--(2,-3)--(0,-3);
\begin{scope}
\clip (-3,-4) rectangle (3.5,2);
\draw[samples=200,domain=-2:4,smooth,variable=\x] plot (\x,{1*((\x)^3)+-3*((\x)^2)+0*(\x)+1});
\end{scope}
\draw (0,1) circle (1pt) (2,-3) circle (1pt);
\end{tikzpicture}
\end{center}
Khẳng định nào sau đây là đúng?
\choice
{\True Hàm số $y=f(x)$ đồng biến trên khoảng $(-\infty; 0)$}
{Hàm số $y=f(x)$ đồng biến trên khoảng $(1;+\infty)$}
{Hàm số $y=f(x)$ nghịch biến trên khoảng $(0; 3)$}
{Hàm số $y=f(x)$ đồng biến trên khoảng $(0; 2)$}
\loigiai{
Dựa vào đồ thị ta thấy hàm số đồng biến trên khoảng $(-\infty; 0)$ và $(2;+\infty)$.
}
\end{ex}

\begin{ex}%[0D3N2-1]
Hàm số nào sau đây là hàm số bậc hai?
\choice
{$y = 2x + 1$}
{$y = -x^4 + x^2 + 1$}
{\True $y = -x^2 + 1$}
{$y = \dfrac{2}{x^2} + x$}
\loigiai{
Hàm số bậc hai là $y = -x^2 + 1$.
}
\end{ex}

\begin{ex}%[0D3H2-2]
Tìm tất cả các giá trị của tham số $m$ để hàm số $y=-4 x^2+4 m x-m^2+2$ nghịch biến trên $(-2 ;+\infty)$.
\choice
{\True $m \leq-4$}
{$m \leq-2$}
{$m<-4$}
{$m<-2$}
\loigiai{
Hàm số $y=-4 x^2+4 m x-m^2+2$ nghịch biến trên khoảng $\left (\dfrac{m}{2};+\infty\right )$ và đồng biến trên khoảng
$\left (-\infty;\dfrac{m}{2}\right )$. Do đó hàm số đã cho nghịch biến trên $(-2 ;+\infty)$ khi và chỉ khi
\[
\dfrac{m}{2} \leq -2 \Leftrightarrow m \leq -4.
\]
}
\end{ex}

\begin{ex}%[H]%[Nguyễn Văn Cường (Cường NV) - Dự án giảng K10 - K11]%[0D6H1-1]
Một hình chữ nhật có các cạnh $x = 4{,}2$ m $\pm 1$ cm, $y = 7$ m $\pm 2$ cm. Chu vi của hình chữ nhật và sai số tuyệt đối của giá trị đó.
\choice
{$22{,}4$ m và $3$ cm}
{$22{,}4$ m và $1$ cm}
{$22{,}4$ m và $2$ cm}
{\True $22{,}4$m và $6$ cm}
\loigiai{
Ta có chu vi hình chữ nhật là $P=2(x+y)=22{,}4$ m $\pm 6$ cm.}
\end{ex}

\begin{ex}%[0-HK1-CT-1-BuiThiXuan-TPHCM-2324]%[VN-MT-6, VM013]%[0D6H4-2]
Số bàn thắng mà một đội bóng ghi được ở mỗi trận đấu trong một mùa giải được thống kê lại ở bảng sau
\begin{center}
\begin{tabular}{|c|c|c|c|c|c|c|}
\hline
Số bàn thắng&$0$&$1$&$2$&$3$&$4$&$5$\\
\hline
Số trận&$4$&$8$&$6$&$3$&$2$&$1$\\
\hline
\end{tabular}
\end{center}
Tìm khoảng tứ phân vị của số bàn thắng trong bảng trên
\choice
{$1$}
{\True $1{,}5$}
{$3$}
{$2$}
\loigiai{
Cỡ mẫu của mẫu số liệu trên là $n=24$.\\
Trung vị (tứ phân vị thứ hai) của mẫu số liệu trên là \[Q_2=\dfrac{x_{12}+x_{13}}{2}=\dfrac{1+2}{2}=1{,}5.\]
Trung vị cho nửa mẫu số liệu bên trái (tứ phân vị thứ nhất) là \[Q_1=\dfrac{x_6+x_7}{2}=\dfrac{1+1}{2}=1.\]
Trung vị cho nửa mẫu số liệu bên phải (tứ phân vị thứ ba) là \[Q_3=\dfrac{x_{18}+x_{19}}{2}=\dfrac{2+3}{2}=2{,}5.\]
Khoảng tứ phân vị của mẫu số liệu là \[\triangle{Q}=Q_3-Q_1=2{,}5-1=1{,}5.\]
}
\end{ex}

\begin{ex}%[0D0N1-2]%[Dự án đề kiểm tra Toán khối 10 HKII NH23-24-Đợt 15-Lê Hải Phụng]%[THPT Chuyên Lê Quý Đôn - HCM]
Gieo một con xúc xắc cân đối và đồng chất hai lần. Xác định số phần tử của không gian mẫu.
\choice
{$ n\left(\Omega\right)=6$}
{\True $ n\left(\Omega\right)=36$}
{$ n\left(\Omega\right)=2$}
{$ n\left(\Omega\right)=12$}
\loigiai{Gieo một con xúc xắc cân đối và đồng chất hai lần thì $ n\left(\Omega\right)=6\cdot6 =36$.}
\end{ex}

\begin{ex}%[0D0N1-3]%[Dự án đề kiểm tra Toán 10 GHKI NH23-24- Lam Nguyen]%[THPT Le Quý Đôn - TPHCM]
Xét phép thử  \lq\lq Gieo một xúc xắc hai lần liên tiếp \rq\rq. Gọi $A$ là biến cố \lq\lq Tổng số chấm hai lần gieo không nhỏ hơn $10$ \rq\rq. Số kết quả thuận lợi cho $A$ là
\choice
{$3$}
{\True $4$}
{$5$}
{$6$}
\loigiai{Số kết quả thuận lợi của biến cố $A$ là các phần tử $(5;5)$, $(6;5)$, $(5;6)$, $(6;6)$.
}
\end{ex}

\begin{ex}%[10-11-12EX-GHK2-2425]%[Dương Quang]%[0H9H3-1]
Trong mặt phẳng $Oxy$, cho đường thẳng $d\colon\heva{&x=2-t \\&y=5+3t}$ ($t\in\mathbb{R}$). Điểm nào sau đây thuộc đường thẳng?
\choice
{$P(2;8)$}
{\True $Q(1;8)$}
{$N(1;5)$}
{$Q(2;-5)$}
\loigiai{
Tại $t=1$ thì $\heva{&x=1 \\&y=8}$. Do đó $Q(1;8)\in d$.
}
\end{ex}

\begin{ex}%[0H9H3-2]%[Dự án đề kiểm tra Toán khối 10 HK2-NH23-24-Đợt 15-Dương Công Tạo]%[THPT Đặng Huy Trứ - Huế]
Cho đường thẳng $d$ có phương trình tham số $\heva{&x=5+2t \\&y=-1-t.&}$ Phương trình tổng quát của đường thẳng qua $B(1; 1)$ và vuông góc với đường thẳng $d$ là
\choice
{$2x-y-3=0$}
{$x+2y-3=0$}
{\True $2x-y-1=0$}
{$2x+y-3=0$}
\loigiai{
Gọi $\Delta$ là đường thẳng cần tìm.\\
Ta có $\Delta \perp d$ suy ra véc-tơ pháp tuyến của $\Delta$ chính là véc-tơ chỉ phương của $d$.\\
Tức là $\vec{n}_{\Delta}=\vec{u}_d=(2;-1)$.\\
Khi đó, phương trình tổng quát của đường thẳng $\Delta$ đi qua $B(1; 1)$ và nhận $\vec{n}_{\Delta}=(2;-1)$ làm véc-tơ pháp tuyến là \[2(x-1)-1(y-1)=0\text{ hay } 2x-y-1=0.\]
}
\end{ex}

\begin{ex}%[0H9H3-4]
Cho đường thẳng $d_{1} \colon 3x+4y+1=0$ và $d_{2}  \colon \heva{& x=15+12t \\& y=1+5t}$. Tính cosin của góc tạo bởi giữa hai đường thẳng đã cho.
\choice{$\dfrac{56}{65}$}
{$-\dfrac{33}{65}$}
{$\dfrac{6}{65}$}
{\True $\dfrac{33}{65}$}
\loigiai{
Ta có $\overrightarrow{n_{1}}=\left(3;4\right) \Rightarrow \overrightarrow{u_{1}}=\left(-4;3\right); \overrightarrow{u_{2}}=\left(12;5\right)$.\\
$\cos \left(d_{1} ;d_{2} \right)=\left|\cos \left(\overrightarrow{u_{1}};\overrightarrow{u_{2}}\right)\right|=\dfrac{\left|-4.12+3.5\right|}{\sqrt{\left(-4\right)^{2} +3^{2}} .\sqrt{12^{2} +5^{2}}} =\dfrac{33}{65}$.}
\end{ex}

\begin{ex}%[0H9N5-1]%[Dự án đề kiểm tra Toán 10 HKII NH23-24- Đoàn Thanh Phong]%[THPT Thiệu Hóa - Thanh Hóa]
Đường Elip $\dfrac{x^2}{16}+\dfrac{y^2}{7}=1$ có tiêu cự bằng
\choice
{$\dfrac{9}{16}$}
{$3$}
{\True $6$}
{$\dfrac{6}{7}$}
\loigiai{
Ta có $a^2=16$, $b^2=7$ nên $c^2=a^2-b^2=9\Rightarrow c=3$.\\
Vậy Elip có tiêu cự bằng $2c=2\cdot3=6$.
}
\end{ex}

\begin{ex}%[0H9N5-5]
Phương trình nào sau đây là phương trình của đường hypebol?
\choice
{\True $\dfrac{x^2}{9} - \dfrac{y^2}{16} = 1$}
{$\dfrac{x^2}{9} + \dfrac{y^2}{16} = 1$}
{$\dfrac{x^2}{9} - \dfrac{y^2}{16} = 0$}
{$\dfrac{x^2}{9} + \dfrac{y^2}{16} = 0$}
\loigiai{
$\dfrac{x^2}{9} - \dfrac{y^2}{16} = 1$ là dạng đúng phương trình của một hypebol.
}
\end{ex}
\TNTF
\begin{ex}%[VN-MT-9]%[0-TK-HK1-KN-1-2425,Nguyễn Hữu Duy]%[0D6V4-2]
Trong một cuộc thi thể thao, người ta ghi lại thời gian (phút) hoàn thành chặng đường đua của một số vận động viên ở bảng sau
\begin{center}
\begin{tabular}{|c|c|c|c|c|c|}
\hline
\textbf{Thời gian (đơn vị: phút)}&$4$&$5$&$6$&$7$&$8$ \\
\hline
\textbf{Số vận động viên}&$3$&$4$&$3$&$5$&$1$ \\
\hline
\end{tabular}
\end{center}
\choiceTF
{\True Mốt của mẫu số liệu trên là $7$ (phút)}
{Tứ phân vị thứ nhất và thứ ba của mẫu số liệu trên lần lượt là $5$ (phút) và $6$ (phút)}
{Độ lệch chuẩn (làm tròn đến hàng phần trăm) là $1{,}23$ (phút)}
{\True Sai số tuyệt đối của số quy tròn độ lệch chuẩn đến hàng phần trăm lớn hơn $0{,}004$ (phút)}

\loigiai{
\begin{itemchoice}
\itemch \textbf{Đúng}.\\
Dữ liệu có giá trị $7$ xuất hiện nhiều nhất ($5$ lần) nên mốt bằng $7$.
\itemch \textbf{Sai}.\\
Cỡ mẫu $n=3+4+3+5+1=16$.\\
Sắp xếp mẫu số liệu theo thứ tự không giảm $4$ $4$ $4$ $5$ $5$ $5$ $5$ $6$ $6$ $6$ $7$ $7$ $7$ $7$ $7$ $8$.\\
Vì $n=16$, là số chẵn nên tứ phân vị thứ hai là $Q_{2}=\dfrac{1}{2}\left(x_{8}+x_{9}\right)=\dfrac{1}{2}(6+6)=6$ (phút).\\
Tứ phân vị thứ nhất $Q_{1}$ là trung vị của mẫu số liệu $4$ $4$ $4$ $5$ $5$ $5$ $5$ $6$.\\
Do đó $Q_{1}=\dfrac{1}{2}(5+5)=5$ (phút).\\
Tứ phân vị thứ ba $Q_{3}$ là trung vị của mẫu số liệu $6$ $6$ $7$ $7$ $7$ $7$ $7$ $8$.\\
Do đó $Q_{3}=\dfrac{1}{2}(7+7)=7$ (phút).
\itemch \textbf{Sai}.\\
Thời gian trung bình của các vận động viên là
\begin{align*}
\overline{x}=\dfrac{4\cdot 3+5\cdot 4+6\cdot 3+7\cdot 5+8\cdot 1}{3+4+3+5+1}=5{,}8125 (\text{phút}).
\end{align*}
Phương sai của mẫu số liệu là
\begin{align*}
s^2=\dfrac{1}{16}\left(3\cdot 4^2+4\cdot 5^2+3\cdot 6^2+5\cdot 7^2+1\cdot 8^2\right)-(5{,}8125)^2=\dfrac{391}{256}.
\end{align*}
Độ lệch chuẩn là
\begin{align*}
s=\sqrt{s^2}=\sqrt{\dfrac{391}{256}} \approx 1{,}24 (\text{phút}).
\end{align*}
\itemch \textbf{Đúng}.\\
Sai số tuyệt đối của số quy tròn độ lệch chuẩn đến hàng phần trăm bằng $\left|\sqrt{\dfrac{391}{256}} - 1{,}24\right| > 0{,}004$ (phút)
\end{itemchoice}
}
\end{ex}

\begin{ex}%[0H9H3-2]  %[Dự án đề kiểm tra Toán 10 HKII NH23-24- Thầy Hải Toán]%[THPT Đống Đa Hà Nội]
Cho hai điểm $A(3;-3)$, $B(-1;-5)$ và đường thẳng $(d)\colon 4x-3y-2=0$
\choiceTF
{Đường thẳng đi qua điểm $A$ và vuông góc với $(d)$ có phương trình $4x+3y=3$}
{Đường tròn đường kính $AB$ có phương trình $(x-1)^2+(y+4)^2=20$}
{Khoảng cách từ $A$ tới $(d)$ nhỏ hơn khoảng cách từ $B$ tới $(d)$}
{\True Cosin của góc tạo bởi $(d)$ và đường thẳng $AB$ bằng $\dfrac{2}{\sqrt{5}}$}
\loigiai{
\begin{enumerate}[a)]
\item Đường thẳng vuông góc với đường thẳng $(d)$ nên có dạng: $\Delta\colon 3x+4y+m=0$.\\
$\Delta$ đi qua $A$ nên thay toạ độ $A$ vào phương trình đường thẳng $\Delta$ ta có \[3\cdot3+4\cdot(-3)+m=0\Leftrightarrow m=3.\]
Vậy phương trình đường thẳng cần tìm là $\Delta\colon 3x+4y+3=0$.
\item Đường tròn đường kính $AB$ thì $AB$ đi qua tâm $I$ và $I$ là trung điểm của $AB$ suy ra $I\left(2;-4\right)$. Suy ra $\vv{IA}=\left(1;1\right)$\\
Bán kính $R=\big|\vv{IA}\big|=\sqrt{1^2+1^2}=\sqrt{2}$.\\
Vậy phương trình đường tròn đường kính $AB$ là $(C)\colon \left(x-2\right)^2+\left(y+4\right)^2=2$.
\item Khoảng cách từ $A$ đến đường thẳng $(d)$ là  $\mathrm{d}\left(A,d\right)=\dfrac{\big|4\cdot 3-3\cdot(-3)-2\big|}{\sqrt{4^2+(-3)^2}}=\dfrac{19}{5}$.\\
Khoảng cách từ $B$ đến đường thẳng $(d)$ là $\mathrm{d}\left(B,d\right)=\dfrac{\big|4\cdot(-1)-3\cdot(-5)-2\big|}{\sqrt{4^2+(-3)^2}}=\dfrac{9}{5}$.\\
Vậy $\mathrm{d}\left(A,d\right)>\mathrm{d}\left(B,d\right)$.
\item Đường thẳng $AB$ nhận $\vv{AB}=\left(-4;-2\right)$ làm một véc-tơ chỉ phương.\\
đường thẳng $(d)$ có véc-tơ pháp tuyến là $n_d=\left(4;-3\right)$ suy ra véc-tơ chỉ phương là $u_d=\left(3;4\right)$.\\
Vậy $\cos\left(d,AB\right)=\dfrac{\big|-4\cdot 3+(-2)\cdot4\big|}{\sqrt{\left(-4\right)^2+\left(-2\right)^2}\cdot\sqrt{3^2+4^2}}=\dfrac{2\sqrt{5}}{5}$.
\end{enumerate}
}
\end{ex}
\TNSA
\begin{ex}%[0D3H2-3]
Đồ thị hàm số $y=5x^2-2x-8$ cắt tia $Ox$ tại điểm có hoành độ bằng bao nhiêu (làm tròn đến hàng phần trăm)?
\shortans{$1{,}48$}
\loigiai{
Ta có $5x^2-2x-8=0\Leftrightarrow \hoac{&x=\dfrac{1+\sqrt{41}}{5} \\&x=\dfrac{1-\sqrt{41}}{5}}$.\\
Vì giao điểm thuộc tia $Ox$ nên hoành độ giao điểm là dương, do đó hoành độ giao điểm là $x=\dfrac{1+\sqrt{41}}{5}\approx 1{,}48$.
}
\end{ex}

\begin{ex}%[10-11-EX-GHK2-2425]%[Lê Toàn]%[0D7V2-7]
Một khung dây thép hình chữ nhật với chiều dài $40$ cm và chiều rộng $30$ cm được uốn lại thành hình chữ nhật mới với kích thước $(40-x)$ cm và $(30+x)$ cm. Tìm điều kiện của $x$ sao cho diện tích hình chữ nhật mà khung tạo thành sau khi uốn tăng lên.
\loigiai{
Điều kiện $\heva{&40-x>0\\&30+x>0}\Leftrightarrow -30<x<40$.\\
Diện tích của hình chữ nhật ban đầu là $40 \cdot 30=1\, 200$ m$^2$.\\
Diện tích của hình chữ nhật sau khi uốn là $(40-x)(30+x)$ cm$^2$.\\
Diện tích sau khi uốn tăng lên khi và chỉ khi
$(40-x)(30+x)-1\, 200>0
\Leftrightarrow -x^2+10x>0
\Leftrightarrow 0<x<10.$\\
Kết hợp với điều kiện, suy ra $0<x<10$.
}
\end{ex}

\begin{ex}%[H]%[BG10-3in1, Phạm Ngọc Trung]%[0D6H3-2]
Người ta điều tra ngẫu nhiên số cân nặng (đơn vị kg) của $20$ học sinh nữ một trường phổ thông, kết quả được ghi trong bảng sau
\begin{center}
\begin{tabular}{|c|c|c|c|c|c|c|c|}
\hline Số cân nặng & $38$ & $40$ & $43$ & $45$ & $48$ & $50$ & \\
\hline Tần số & $3$ & $3$ & $4$ & $5$ & $2$ & $3$ & $N=20$ \\
\hline
\end{tabular}
\end{center}
Tính số cân nặng trung bình $\overline{x}$ (\textit{kết quả làm tròn đến hàng phần chục}).
\shortans{43{,}9}
\loigiai{Số cân nặng trung bình của $20$ học sinh là
\begin{align*}
\overline{x}=\dfrac{38\cdot 3+40\cdot 3+43\cdot 4+45\cdot 5+48\cdot 2+50\cdot 3}{20}=\dfrac{877}{20}\approx 43{,}9.
\end{align*}
}
\end{ex}

\begin{ex}%[dự án HSA, Huỳnh Xuân Tín]%[0H9V4-2]
Đường tròn $(C)$ có tâm $I$ thuộc đường thẳng $d\colon x+3y+8=0$, đi qua điểm $A\left(-2;1\right)$ và tiếp xúc với đường thẳng $\Delta:3x-4y+10=0$. Tính bán kính đường tròn $(C)$.
\shortans[1]{$5$}
\loigiai{
Ta có $I\in d \Rightarrow I(-3m-8;m)$.
\begin{eqnarray*}
&&AI^2=\left[\mathrm{d}(I,\Delta)\right]^2 \text{ (cùng bằng } R^2) \\
& \Leftrightarrow & \left(-3m-6\right)^2+\left(m-1\right)^2=\dfrac{\left[3(-3m-8)-4m+10\right]^2}{3^2+(-4)^2}\\
& \Leftrightarrow & 10m^2+34m+37=\dfrac{169m^2+364m+196}{25}\\
& \Leftrightarrow & 81m^2+486m+729=0\\
& \Leftrightarrow & m=-3.
\end{eqnarray*}
Do đó, đường tròn $(C)$ có $\heva{&\text{tâm } I(1;-3)\\& \text{bán kính } R=AI=\sqrt{(9-6)^2+(-3-1)^2}=5.}$\\
Vậy phương trình đường tròn là ${\left(x-1\right)}^2+{\left(y+3\right)}^2=25$.
}
\end{ex}
\TL
\begin{ex}%[Mức 1]%[Dự án Giảng 10-11 Nhóm Toán & LaTex, Lê Minh Thiện Anh]%[0D3H1-2]
Tìm tập xác định của hàm số $y=\sqrt{\dfrac{x^2-2x+3}{4-x^2}}$.
\loigiai{
Vì $x^2-2x+3 =(x-1)^2+2>0,\;\forall x \in \mathbb{R}$ nên
hàm số xác định khi và chỉ khi \[\dfrac{x^2-2x+3}{4-x^2} \ge 0\Leftrightarrow 4-x^2>0 \Leftrightarrow -2<x<2.\]
Vậy tập xác định của hàm số là $\mathscr{D}=(-2;2)$.
}
\end{ex}

\begin{ex}%[0D7V3-6]
\immini{Mặt cắt đứng của cột cây số trên quốc lộ có dạng nửa hình tròn ở phía trên và phía dưới có dạng hình chữ nhật (xem hình bên). Biết rằng đường kính của nửa hình tròn cũng là cạnh phía trên của hình chữ nhật và đường chéo của hình chữ nhật có độ dài $66$ cm. Tìm diện tích của hình chữ nhật, biết rằng diện tích của phần nửa hình tròn bằng $0{,}3$ lần diện tích của phần hình chữ nhật. Lấy $\pi=3{,}14$ và làm tròn kết quả đến hàng đơn vị.}{\begin{tikzpicture}[scale=.8, >=stealth]
\coordinate (D) at (2,0);
\coordinate (B) at (-2,-5);
\coordinate (A) at (-2,0);
\coordinate (C) at (2,-5);
\draw[smooth](A)--(B)--(C)--(D)--(A) (B)--(D);
\clip (-3,-7) rectangle (3,2);
\draw[thick,smooth,samples=200,domain
=-2:2] plot(\x,{sqrt(4-(\x)^(2))});
\draw ($(A)!.5!(D)$) node[above]{$x$};
\draw (-.5,-2) node[below,rotate=52.5]{$66$ cm};
\draw ($(B)!.5!(C)$) node[below]{Mặt cắt của cột cây số};
\end{tikzpicture}}
\loigiai{Gọi đường kính của nửa hình tròn là $x~(\mathrm{cm})$ $(x>0)$. \\Độ dài cạnh bên của hình chữ nhật là $\sqrt{66^2-x^2}$. \\Diện tích nửa hình tròn là $\dfrac{3{,}14 x^2}{8}$. \\Diện tích hình chữ nhật là $x \sqrt{66^2-x^2}$. Theo giả thiết ta có
\allowdisplaybreaks
\begin{eqnarray*}
\dfrac{3{,}14 x^2}{8}=0{,}3 x \sqrt{66^2-x^2} &\Leftrightarrow& 157 x=120 \sqrt{66^2-x^2} \\&\Leftrightarrow& x^2=\dfrac{62726400}{39049} \Leftrightarrow x \approx \pm 40{,}08.
\end{eqnarray*}
Vậy diện tích  của hình chữ nhật là \[40{,}08 \cdot 52{,}44\approx 2\,101{,}7952~\left(\mathrm{cm}^2\right).\]
}
\end{ex}

\begin{ex}%[Dự án C THPTQG 2025]%[Vương Quốc Phong]%[0D0C2-9]
Hai bạn Nga và Nhung chơi trò tung xúc xắc. Mỗi bạn tung $1$ con xúc xắc $3$ lần, ai có tổng số chấm $3$ lần gieo lớn hơn thì thắng. Nga chơi trước và được $14$ chấm. Tính xác suất để Nhung thắng Nga.
\loigiai{
Gọi $A$ là biến cố: \lq\lq Nhung thắng Nga sau ba lần tung\rq\rq.\\
Khi đó $\mathrm{P}(A)$ là xác suất tổng số chấm Nga tung được sau ba lần tung lớn hơn $14$. \\
Ta có $\Omega=\left\{(a_1, a_2, a_3)|a_i \in \{1,2,\ldots, 6\}, i=\overline{1,3}\right\}\Rightarrow n(\Omega)=6^3=216$.\\
$A=\left\{(a_1, a_2, a_3)|a_1+a_2+a_3 \ge 15, a_i \in \{1,2,\ldots, 6\}, i=\overline{1,3}\right\}$.\\
Để đếm số phần tử của $A$, ta chia thành các trường hợp sau:
\begin{itemize}
\item Trường hợp 1: $a_1+a_2+a_3=15$, gồm bộ $(5,5,5)$ và các bộ là hoán vị của $(4;5;6)$ và $(3;6;6)$.\\
Trường hợp này có $1+6+3=10$ (bộ)
\item  Trường hợp 2: $a_1+a_2+a_3=16$, gồm các hoán vị của $(5,5,6)$ và $(4,6,6)$.\\ Trường hợp này có $3+3=6$ (bộ)
\item  Trường hợp 3: $a_1+a_2+a_3=17$, gồm các hoán vị của $(5,6,6)$.\\
Trường hợp này có $3$ (bộ).
\item Trường hợp 4: $a_1+a_2+a_3=18$, gồm $(6,6,6)$.\\
Trường hợp này có $1$ (bộ).
\end{itemize}
Vậy $n(A)=20$. \\
Xác suất cần tìm là $\mathrm{P}(A)=\dfrac{n(A)}{n(\Omega)}=\dfrac{20}{216}=\dfrac{5}{54}$.
}
\end{ex}
\Closesolutionfile{ans}
\Closesolutionfile{ansbook}